\chapter*{Conclusiones}
\addcontentsline{toc}{chapter}{Conclusiones}
\markboth{Conclusiones}{Conclusiones}

A partir del análisis técnico, operativo, económico y ambiental desarrollado en esta monografía sobre la energía de biomasa en plantas de potencia, se establecen las siguientes conclusiones:

\begin{enumerate}[leftmargin=1.5em]
\item \textbf{Recurso energético almacenado y condicionantes fisicoquímicas:} La biomasa constituye energía solar almacenada químicamente mediante fotosíntesis en un ciclo corto del carbono. Sus propiedades determinantes ---humedad, contenido de materia volátil (70--85~\%), química de cenizas alcalinas y baja densidad aparente--- definen tanto el diseño de los equipos térmicos (limitando la temperatura de vapor a 450--540~\textdegree C por corrosión) como el radio económico de abastecimiento logístico.

\item \textbf{Adecuación de las rutas de conversión:} No existe una tecnología universalmente superior; la selección depende de la humedad y escala del recurso. Las rutas termoquímicas se destinan a biomasa seca y lignocelulósica, dominando la combustión directa con ciclo Rankine en mediana y gran escala (20--35~\% de eficiencia eléctrica neta dedicada y 36--44~\% en co-combustión), mientras que la gasificación resulta atractiva para sistemas descentralizados con motores y la pirólisis densifica el recurso en bioaceite. Para sustratos muy húmedos y biodegradables, la digestión anaeróbica con motogeneradores (35--42~\% de eficiencia sobre biogás) permite tratar residuos y generar energía simultáneamente.

\item \textbf{Escalas y cogeneración:} Las centrales abarcan desde biodigestores rurales de decenas de kilovatios hasta reconversiones de centrales de carbón de gigavatios como Drax. La cogeneración (CHP) representa la configuración más eficiente (70--90~\% de eficiencia global), lo que explica la estrecha integración histórica de la biomasa con ingenios azucareros, papeleras y aserraderos, ejemplificada regionalmente en Guatemala y a gran escala en Alholmens Kraft.

\item \textbf{Valor estratégico en el sistema de potencia:} El principal aporte de la biomasa frente a la solar y la eólica no radica en el costo de la energía ---cuyo LCOE promedio mundial se sitúa en $0{,}087~\text{USD/kWh}$ y depende fuertemente del combustible---, sino en su carácter despachable (factor de planta medio de $\sim 53~\%$) y en los servicios auxiliares que brindan sus generadores síncronos: inercia rotacional natural, regulación primaria de frecuencia, control de tensión en redes rurales débiles y corriente de cortocircuito.

\item \textbf{Sostenibilidad y perspectivas en Panamá:} La neutralidad climática de la biomasa exige análisis de ciclo de vida que prioricen residuos locales y eviten cambios adversos en el uso del suelo, abriendo además la vía a emisiones netas negativas mediante BECCS. En Panamá, aunque las estadísticas de oferta eléctrica de la ASEP en 2025 no registran centrales dedicadas en red, el país dispone de marco legal (Ley 42 de 2011) y recursos con potencial real ---bagazo de caña, residuos arroceros, efluentes agropecuarios, paja canalera y biogás de rellenos sanitarios--- cuya integración mediante esquemas de cogeneración y generación distribuida puede fortalecer la seguridad y estabilidad de la matriz energética nacional.
\end{enumerate}
